\documentclass[11pt]{article}
\usepackage[margin=1in]{geometry}
\usepackage{amsmath,amssymb}
\usepackage{booktabs}
\usepackage[hidelinks]{hyperref}
\usepackage{url}

\title{A Publicly Verifiable G63 Max-Cut Solution of 27,050\\Found with Population Annealing on a Consumer GPU}
\author{Filip Majchrzak\\
\small SAPIO MORPHEUS / COMPLEX (independent research project), Poland\\
\small \url{https://github.com/kosza123}}
\date{November 2026}

\begin{document}
\maketitle

\begin{abstract}
We report a Max-Cut configuration of value 27,050 for the G-set instance G63 (7,000 vertices, 41,459 unit-weight edges),
improving upon the published reference value of 27,047. The configuration was found with population annealing Monte Carlo,
a known method, implemented for a single consumer GPU (NVIDIA GeForce RTX 3060 Ti), 158.5 minutes into a seven-hour experiment.
The complete partition, checksums and dependency-free verifiers are publicly available, so the value can be checked
independently in under a second. Relative to the published 27,047 configuration, the new solution differs by a single connected
set of 1,092 vertices whose flip adds three cut edges, together with zero-cost clusters. We make no claim of optimality or
of a new algorithm, and we discuss the limitations of the experiment.
\end{abstract}

\section{Introduction}
The G-set graphs, a collection whose first instances (G1--G54) were generated with the \emph{rudy} generator and introduced by
Helmberg and Rendl~\cite{helmberg2000}, later extended with larger instances including G63~\cite{gset}, are a standard benchmark for
Max-Cut heuristics~\cite{benlic2013,shylo2015}. Instance G63 has 7,000 vertices and 41,459 edges with
unit weights; it is one of the two G-set instances with the largest number of edges. For about a decade the best reported cut for G63
was 27,045, attributed to global equilibrium search~\cite{shylo2015}. In 2025, Khan and Shukla reported 27,047 with a GPU-accelerated
population annealing framework and published the corresponding configuration~\cite{khan2025}; a later preprint by the same group again
lists 27,047 as the best value it achieved for G63~\cite{khan2026}.

Here we present a configuration of value 27,050. Our contribution is computational: a verifiable configuration that exceeds the
published reference, obtained with a standard method on consumer hardware, together with the data needed to check it.

\section{The configuration and its verification}
The partition is given as a string of 7,000 characters \texttt{0}/\texttt{1}, character $k$ being the side of vertex $k+1$ in the
G-set numbering, and equivalently as 1,750 hexadecimal digits, each read most significant bit first (the hexadecimal format of~\cite{khan2025};
reading their string this way reproduces their value of 27,047).
Its SHA-256 digest is \texttt{2e3d330e\ldots a566df}; the original G63 file~\cite{gset} has digest \texttt{a1d08e1e\ldots a99b}
(full values in the repository).

The cut value was recomputed by seven independent implementations (Python, Node.js, awk, several numpy formulations including the spin form
$(\sum w - x^{\top}Wx)/2$ and the Laplacian form $x^{\top}Lx/4$, PowerShell and gawk), three of them written by separate AI-agent instances
without access to the search code and one run on a freshly downloaded copy of G63. All give 27,050; the complementary partition
gives the same value. As controls, the published configuration of~\cite{khan2025} gives 27,047 with the same tools, and a
one-position shift of our string gives 20,744 or 20,588, close to the value $m/2 = 20{,}729.5$ expected for a random partition.
The configuration is a 1-flip local optimum (no single-vertex flip increases the cut; 221 vertices have zero gain).

\paragraph{Relation to the published configuration.}
Up to global complement, the two configurations differ on 1,140 vertices, which form 43 connected components in G63.
Exactly one component, a connected set of 1,092 vertices, changes the cut: flipping it in the published configuration adds
three edges, giving 27,050. The other 42 components are zero-cost clusters. The improvement is therefore a coordinated
rearrangement of about one thousand vertices rather than a small local modification.

\section{Method}
Population annealing~\cite{hukushima2003,machta2010,wang2015} evolves a population of replicas while the inverse temperature
$\beta$ increases; at each step the population is resampled in proportion to Boltzmann weights and every replica is updated by
Monte Carlo sweeps. We maximise the cut with weights $\propto \exp(\beta\,\mathrm{cut})$.

The run that produced the configuration used $R=8{,}192$ replicas initialised uniformly at random, a geometric schedule of
6,000 steps from $\beta=0.3$ to $\beta=9.0$, systematic resampling with weights $\exp((\beta_t-\beta_{t-1})\,\mathrm{cut})$,
and ten heat-bath sweeps per step, with flip probability $1/(1+e^{-\beta g})$ for a cut change $g$. Vertices were updated by the
classes of a greedy proper colouring with eight colours, so that a whole class is updated in parallel; every anneal ended with a
zero-temperature quench. The code was written for this experiment in CUDA C, compiled at run time through CuPy, with spins stored
in a vertex-major, replica-minor layout for coalesced memory access; cuts were tracked incrementally and periodically recounted.
One anneal took about ten minutes on an RTX 3060 Ti that was shared with a second annealing process.

\section{Results}
The experiment started on 7 November 2026 at 23:35:21 CET with a hard limit of seven hours. A first configuration above the
reference (27,048) appeared after 113 minutes, and the 27,050 configuration after 158.5 minutes (anneal 12 of the run described
above). No configuration above 27,050 was found in the remaining time.

Table~\ref{tab:stats} summarises the best cut of 66 complete anneals with population 8,192 and three schedules chosen during the
experiment ($\beta$ up to 7, 9 and 10; 4,000, 6,000 and 9,000 steps). For 55 of these anneals (those whose saved states included a 27,047 state) we checked one 27,047 state; all
of them belonged to the family of the published configuration (differing from it only by zero-cost clusters). Shorter anneals
(population 2,048--4,096, at most 1,500 steps) ended between 26,957 and 27,045.

\begin{table}[h]
\centering
\begin{tabular}{cc}
\toprule
best cut of the anneal & number of anneals \\
\midrule
27,046 & 2 \\
27,047 & 54 \\
27,048 & 9 \\
27,050 & 1 \\
\bottomrule
\end{tabular}
\caption{Best cut per anneal for 66 complete anneals with population 8,192 started from random states.}
\label{tab:stats}
\end{table}

Before population annealing, tabu search and exact re-optimisation of ball-shaped regions of 300--1,200 vertices by
mixed-integer programming around the published configuration found no improvement. The 27,047 level behaves as a large plateau:
about 235 vertices of a 27,047 configuration have zero flip gain, and different 27,047 configurations differ by independent
zero-cost clusters.

\section{Discussion and limitations}
The result is a single verified configuration on a single instance, obtained in one experiment whose parameters were adjusted
during the run; we did not compare solvers on equal hardware. Population annealing is a known method, and the run times reported
in~\cite{khan2025} (runs that produced their best cut typically ran 8--24 hours on an RTX A6000) are run lengths under different
settings, not comparable with our anneal times; no equal-hardware or equal-time comparison was made. We make no claim that 27,050 is optimal.

All primary sources we could verify list at most 27,047 for G63~\cite{khan2025,khan2026}. A secondary aggregator page has been
described in search-engine snippets as mentioning a value of 27,052; we could not access that page or locate a primary source or
configuration. If such a configuration exists, it exceeds the value reported here.

\section*{Data availability}
The partition, checksums, verifiers, verification log and methodology are available at
\url{https://github.com/kosza123/g63-maxcut-27050}.

\section*{Use of AI}
The experiment was conceived and directed by the author. The search code, the execution of the experiment, the verification
tools and a first draft of this manuscript were produced with an AI assistant (Claude, Anthropic) acting as a research operator
under the author's direction; three of the independent recounts were written by separate AI-agent instances without access to the
search code. The author is responsible for the content.

\begin{thebibliography}{9}
\bibitem{helmberg2000} C. Helmberg and F. Rendl, A spectral bundle method for semidefinite programming, \emph{SIAM Journal on Optimization} 10(3):673--696, 2000.
\bibitem{gset} G-set Max-Cut test graphs, collection hosted by Y. Ye, Stanford University, \url{https://web.stanford.edu/~yyye/yyye/Gset/} (accessed November 2026).
\bibitem{benlic2013} U. Benlic and J.-K. Hao, Breakout local search for the max-cut problem, \emph{Engineering Applications of Artificial Intelligence} 26(3):1162--1173, 2013.
\bibitem{shylo2015} V. P. Shylo, F. Glover and I. V. Sergienko, Teams of global equilibrium search algorithms for solving the weighted maximum cut problem in parallel, \emph{Cybernetics and Systems Analysis} 51(1):16--24, 2015.
\bibitem{hukushima2003} K. Hukushima and Y. Iba, Population annealing and its application to a spin glass, \emph{AIP Conference Proceedings} 690:200--206, 2003.
\bibitem{machta2010} J. Machta, Population annealing with weighted averages: A Monte Carlo method for rough free-energy landscapes, \emph{Physical Review E} 82:026704, 2010.
\bibitem{wang2015} W. Wang, J. Machta and H. G. Katzgraber, Population annealing: Theory and application in spin glasses, \emph{Physical Review E} 92:063307, 2015.
\bibitem{khan2025} N. Khan and N. Shukla, New best-known Max-Cut solution for the G63 instance in the G-set benchmark, arXiv:2510.21105, 2025.
\bibitem{khan2026} N. Khan, R. Redding and N. Shukla, Leveraging population dynamics to steer efficient search in large-scale combinatorial optimization, arXiv:2606.25203, 2026.
\end{thebibliography}
\end{document}
